% Figure: pressure-altitude profile of the International Standard Atmosphere
% to 32 km, compared with the simple isothermal exponential model. The
% isothermal model overshoots aloft because it ignores the cold troposphere.
% Data from data/isa_profile.csv conceptually; curves computed inline. No
% em-dashes.
\begin{figure}[htbp]
\centering
\begin{tikzpicture}
\begin{axis}[
  width=11cm, height=8cm,
  xlabel={pressure $p$ (kPa)},
  ylabel={geometric altitude $z$ (km)},
  xmin=0, xmax=105, ymin=0, ymax=32,
  axis lines=left,
  legend pos=north east,
  legend cell align=left,
  grid=major, grid style={enggray!20},
  tick label style={font=\scriptsize},
  label style={font=\small},
  legend style={font=\scriptsize},
]
% ISA approximation, sampled points (troposphere + stratosphere)
\addplot[engblue, very thick, mark=*, mark size=1pt] coordinates {
  (101.3,0) (89.9,1) (79.5,2) (70.1,3) (61.6,4) (54.0,5)
  (47.2,6) (41.1,7) (35.6,8) (30.7,9) (26.5,10) (22.6,11)
  (16.5,13) (12.0,15) (8.8,17) (6.5,19) (5.5,20) (4.0,22)
  (2.5,25) (1.6,28) (1.0,30) (0.87,32)
};
\addlegendentry{ISA (lapse-rate model)}
% isothermal model p = 101.3 * exp(-z/8.4)  with z in km
\addplot[engred, dashed, very thick, domain=0:32, samples=60] {101.3*exp(-x/8.4)};
\addlegendentry{isothermal, $H=8.4$ km}
% tropopause marker
\draw[engamber, dotted, thick] (axis cs:0,11) -- (axis cs:105,11);
\node[engamber, font=\scriptsize, anchor=west] at (axis cs:55,11.8) {tropopause $11$ km};
\end{axis}
\end{tikzpicture}
\caption[Standard-atmosphere pressure profile]{Pressure versus altitude
for the International Standard Atmosphere (blue) compared with the
single-scale-height isothermal model $p=p_0e^{-z/H}$, $H=8.4\,$km (red
dashed). The two agree near the ground, where the scale height is set by
the surface temperature, and diverge aloft: the real troposphere cools
at $6.5\,$K/km, so the air column is denser than the isothermal model
assumes and the pressure falls faster than a single exponential. Above
the tropopause (amber, $11\,$km) the stratosphere is isothermal and the
two curves run parallel again on a log scale. Altimeters are calibrated
to the ISA, so departures of the real atmosphere from it appear directly
as altitude error. ISA values from the lapse-rate integration in
\texttt{code/standard\_atmosphere.py}.}
\label{fig:vol03ch10:atmosphere}
\end{figure}
